% TOP_PROBLEM: {"id": 30007201, "title": "Coefficient injectivity for unimodular one-variable equations", "category_id": 17, "category": {"id": 17, "name": "group_theory", "display_name": "Group Theory", "description": "Problems about groups, group actions, representations, and related algebraic structures.", "slug": "group-theory", "order_index": 17, "created_at": "2026-07-31T15:26:25.671Z"}, "status": "open"}
% FIRST_POSTED: 2026-10-06
% ATTEMPT_STATUS: solved
% OPENAI_CLAIM: {"schema_version":1,"source_repo":"https://github.com/openai/math","source_commit":"adc7f1241b42e322a6451854ab7e4b4c146bf78a","release_date":"2026-10-06","match":"full","resolution":"proved","independently_reviewed":false,"second_pass":"agreed","families":[{"family":"256","title":"Nonsingular systems of equations over arbitrary groups","subject":"Group theory","manuscripts":[{"title":"The Kervaire theorem for groups","pdf_path":"preprints/The-Kervaire-Theorem-for-Groups-September-24-2026/The-Kervaire-Theorem-for-Groups-September-24-2026.pdf","readme_path":"preprints/The-Kervaire-Theorem-for-Groups-September-24-2026/README.md","date":"2026-09-24","theorem_ref":"preprints/The-Kervaire-Theorem-for-Groups-September-24-2026/build/sections/introduction.tex:23, result labelled thm:main (pinned adc7f1241b42e322a6451854ab7e4b4c146bf78a)"}],"lean":{"doc_path":"lean/docs/256.md","comparators":["lean/ComparatorChallenges/Kervaire.json","lean/ComparatorChallenges/Kervaire.lean"],"declarations":[],"covers_main_theorem":true},"reasoning_trace":null}]}
\documentclass[11pt]{article}
\usepackage[margin=25mm]{geometry}
\usepackage{fontspec}
\setmainfont{Latin Modern Roman}
\usepackage{amsmath,amssymb,mathtools}
\usepackage{unicode-math}
\setmathfont{Latin Modern Math}
\usepackage{enumitem}
\usepackage{xurl,hyperref,bookmark}
\hypersetup{colorlinks=true,urlcolor=blue}
\setcounter{secnumdepth}{0}
\setlength{\parindent}{0pt}
\setlength{\parskip}{5pt}
\setlength{\emergencystretch}{3em}
\begin{document}
\section*{Coefficient injectivity for unimodular one-variable equations}

\subsection{OpenAI's claim}
OpenAI presents a proof in the selected manuscript. For every nontrivial group $G$, the free product $G*\mathbb Z$ is claimed not normally generated by a single element.

\subsection{Outline of the argument}
A traced spectral phase provides a unitary with a controlled large fixed space. Translating a putative coefficient relation into a planar surface yields one nontrivial boundary, which the phase calculation excludes; this prevents a normal generator of the required form.

\subsection{Scope relative to this problem}
The Kervaire paper main theorem explicitly allows every coefficient group A and every unimodular w in A*<t>, meaning exponent sum plus/minus one, and proves the natural coefficient homomorphism injective. lean/docs/256.md documents exactly this selected theorem. Within these assumptions, this is a claimed affirmative answer to the stated question.

\subsection{Formal verification}
For family 256, the selected statement is coefficient injectivity for a single relator of exponent sum $1$ or $-1$ in $A*\mathbb Z$. It underlies the Kervaire consequence but does not include the newer general nonsingular-system theorem. The listed Comparator files are \texttt{Kervaire.lean}. The inventory does not supply a named declaration for these Comparator entries. This site has not executed the Lean code. The scope descriptions do not certify the remaining manuscript arguments.

\subsection{Status caveat}
OpenAI attributes these manuscripts to work by an internal model. Its release cautions that arguments without formalization can contain mistakes. This catalogue has not independently checked the proofs; the labels report the mathematical claims attributed to OpenAI.

\subsection{Sources}
\begin{itemize}
\item OpenAI, \emph{The Kervaire theorem for groups}, 2026-09-24; citation key \texttt{OAI:The-Kervaire-Theorem-for-Groups-September-24-2026}. \href{https://github.com/openai/math/blob/adc7f1241b42e322a6451854ab7e4b4c146bf78a/preprints/The-Kervaire-Theorem-for-Groups-September-24-2026/The-Kervaire-Theorem-for-Groups-September-24-2026.pdf}{PDF (GitHub)}; \href{https://github.com/openai/math/raw/adc7f1241b42e322a6451854ab7e4b4c146bf78a/preprints/The-Kervaire-Theorem-for-Groups-September-24-2026/The-Kervaire-Theorem-for-Groups-September-24-2026.pdf}{PDF (direct)}; \href{https://github.com/openai/math/blob/adc7f1241b42e322a6451854ab7e4b4c146bf78a/preprints/The-Kervaire-Theorem-for-Groups-September-24-2026/README.md}{Manuscript page and citation}; \href{https://github.com/openai/math/blob/main/preprints/The-Kervaire-Theorem-for-Groups-September-24-2026/The-Kervaire-Theorem-for-Groups-September-24-2026.pdf}{Latest version}.
\item Family 256: \href{https://github.com/openai/math/blob/adc7f1241b42e322a6451854ab7e4b4c146bf78a/lean/docs/256.md}{Lean scope}.
\item \href{https://github.com/openai/math/blob/adc7f1241b42e322a6451854ab7e4b4c146bf78a/lean/ComparatorChallenges/Kervaire.json}{Comparator Kervaire.json}.
\item \href{https://github.com/openai/math/blob/adc7f1241b42e322a6451854ab7e4b4c146bf78a/lean/ComparatorChallenges/Kervaire.lean}{Comparator Kervaire.lean}.
\item \href{https://github.com/openai/math/blob/adc7f1241b42e322a6451854ab7e4b4c146bf78a/CONTENTS.md}{OpenAI catalogue entry}.
\end{itemize}
\end{document}
